% ---------------------------------------------------------------------
% Appendix C
% ---------------------------------------------------------------------

\begin{cnread}
本附录给出第 9 节用到的库仑径向作用量的完整初等推导，并对三个参数点做
数值核对（V7.1）。结果是半作用量
$\int_{r_-}^{r_+}p_r\,\mathrm dr=\pi\bigl(m\alpha_e/\sqrt{-2mE}-l\bigr)$，
整体作用量是它的两倍。由此立刻得到玻尔能谱。
\end{cnread}

\subsection{The radial momentum and the turning points}

For a Coulomb pair with potential $V=-k/r$ and conserved angular momentum $l$,
\begin{equation}\label{eq:pr}
  p_r^{2}=2m\Bigl(E+\frac kr\Bigr)-\frac{l^{2}}{r^{2}}
  =\frac{2mE\,r^{2}+2mk\,r-l^{2}}{r^{2}} .
\end{equation}
For $E<0$ put $A=2m|E|>0$, so that $2mE=-A$.  The turning points are the positive
roots of $Ar^{2}-2mk\,r+l^{2}=0$:
\begin{equation}\label{eq:turning}
  r_{\pm}=\frac{mk\pm\sqrt{m^{2}k^{2}-Al^{2}}}{A},
  \qquad
  r_++r_-=\frac{2mk}{A},
  \qquad
  r_+r_-=\frac{l^{2}}{A} .
\end{equation}
Between them $p_r=\sqrt A\,\sqrt{(r_+-r)(r-r_-)}/r$.

\subsection{The half-action}

\begin{lemma}[Half radial action]
\label{lem:half-action}
With $A=2m|E|$ and the turning points \eqref{eq:turning},
\begin{equation}\label{eq:half-action}
  \int_{r_-}^{r_+}p_r\,\mathrm dr
  =\pi\Bigl(\frac{mk}{\sqrt{A}}-l\Bigr)
  =\pi\Bigl(\frac{mk}{\sqrt{-2mE}}-l\Bigr).
\end{equation}
\end{lemma}

\begin{proof}
Substituting $p_r=\sqrt A\sqrt{(r_+-r)(r-r_-)}/r$ reduces the claim to the
elementary integral
\begin{equation}\label{eq:elementary}
  \int_{a}^{b}\frac{\sqrt{(b-x)(x-a)}}{x}\,\mathrm dx
  =\frac{\pi(a+b)}{2}-\pi\sqrt{ab},
\end{equation}
which is Check~V7.4 at three parameter points.  Applying it with $a=r_-$,
$b=r_+$ and using \eqref{eq:turning},
\[
  \int_{r_-}^{r_+}p_r\,\mathrm dr
  =\sqrt A\,\pi\Bigl(\frac{r_++r_-}{2}-\sqrt{r_+r_-}\Bigr)
  =\sqrt A\,\pi\Bigl(\frac{mk}{A}-\frac{l}{\sqrt A}\Bigr)
  =\pi\Bigl(\frac{mk}{\sqrt A}-l\Bigr).
\]
The second form uses $A=-2mE$.
\end{proof}

\begin{remark}[On \eqref{eq:elementary}]
The integral is a standard contour integral; it can also be obtained by the
Kepler substitution $x=\bar r(1-e\cos\psi)$ with $\bar r=(a+b)/2$ and
$e=(b-a)/(b+a)$.  The version stated is the one verified numerically.
\end{remark}

\subsection{Numerical verification}

Check~V7.1 evaluates the left side of \eqref{eq:half-action} by symbolic
quadrature and compares with the closed form at three parameter points
$(m,k,l,E)$:

\begin{center}\small
\renewcommand{\arraystretch}{1.2}
\begin{tabular}{@{}ccccc@{}}
\toprule
$m$ & $k$ & $l$ & $E$ & $\int_{r_-}^{r_+}p_r\,\mathrm dr$ \\
\midrule
$1$ & $1$ & $1/2$ & $-1/2$ & $1.57079632679\ldots=\pi/2$ \\
$2$ & $3$ & $1$ & $-9/8$ & $5.74417322273\ldots$ \\
$1/2$ & $3$ & $1$ & $-8/9$ & $1.85665065184\ldots$ \\
\bottomrule
\end{tabular}
\end{center}

In each case the quadrature agrees with
$\pi\bigl(mk/\sqrt{-2mE}-l\bigr)$ to better than $10^{-12}$.

\subsection{The spectrum}

With $\alpha_e=e^{2}/(4\pi)$ and $k=\alpha_e$, the full action is twice
\eqref{eq:half-action}, and Quantization~\eqref{eq:bohr-sommerfeld} gives
\begin{equation}\label{eq:quantization}
  2\pi\Bigl(\frac{m\alpha_e}{\sqrt{-2mE}}-l\Bigr)=2\pi n_r\hbar
  \quad\Longleftrightarrow\quad
  \frac{m\alpha_e}{\sqrt{-2mE}}=l+n_r\hbar=\hbar n ,
\end{equation}
where $n:=n_r+l/\hbar$.  Then
\begin{equation}\label{eq:spectrum-app}
  E_n=-\frac{m\alpha_e^{2}}{2\hbar^{2}n^{2}} .
\end{equation}
Check~V7.2 substitutes \eqref{eq:spectrum-app} into the middle expression of
\eqref{eq:quantization} and finds $\hbar n$ identically; Check~V7.3 recovers
$E_1=-m\alpha_e^{2}/(2\hbar^{2})$.

\begin{warning}[Semiclassical status]
As stated in Warning~\ref{warn:semiclassical}, \eqref{eq:spectrum-app} is a
Bohr--Sommerfeld result.  Its agreement with the exact hydrogen spectrum is a
property of the Coulomb potential; no operator-theoretic statement is made, and
the multiplicities $n^{2}$ of the levels are not derived.
\end{warning}
